% Kurzfassung (DE) und Abstract (EN) -- beide Sprachen werden unabhängig von
% \thesislanguage ausgegeben; die Hauptsprache zuerst.
% Richtwert: je 200--300 Wörter. Struktur: Problem -> Ansatz -> Daten/Methode ->
% Hauptergebnis (mit Zahl) -> Bedeutung.

\newcommand{\kurzfassungDE}{%
\begin{otherlanguage}{ngerman}
\chapter*{Kurzfassung}
\todo{Kurzfassung schreiben.} Die Röntgen-Nahfeldholographie an
Synchrotronquellen wie PETRA~III rekonstruiert komplexe Objekte aus
Inline-Hologrammen. Die Rekonstruktion setzt eine genaue Kenntnis der
Aufnahmegeometrie voraus, die in der Fresnel-Zahl $\Fr$ zusammengefasst wird;
deren Bestimmung ist das Autofokus-Problem. Modellbasierte Verfahren lösen es
durch Optimierung, sind aber langsam. Diese Arbeit untersucht, ob ein
neuronales Netz $\Fr$ bzw. den Abstand $\zOne$ direkt aus dem Hologramm
schätzen kann. \todo{Daten, Modell, Hauptergebnis mit Zahl, Einordnung.}
\end{otherlanguage}
}

\newcommand{\abstractEN}{%
\begin{otherlanguage}{english}
\chapter*{Abstract}
\todo{Write abstract.} X-ray near-field holography at synchrotron sources such
as PETRA~III reconstructs complex-valued objects from in-line holograms. The
reconstruction requires accurate knowledge of the acquisition geometry, which is
summarised by the Fresnel number $\Fr$; determining it is the autofocus problem.
Model-based methods solve it by optimisation but are slow. This thesis
investigates whether a neural network can estimate $\Fr$ or the distance
$\zOne$ directly from a hologram. \todo{Data, model, main result with a number,
significance.}
\end{otherlanguage}
}

\ifdefstring{\thesislanguage}{german}{\kurzfassungDE\abstractEN}{\abstractEN\kurzfassungDE}
